\section*{Bab \ref{ch:g2:where-animals-live} --- Tempat Hewan Hidup}

\begin{solution}{exo:g2:where-animals-live:1}
Jenis tempat yang secara alami dihuni seekor \omterm{def:g1:animals-around-us:animal}{hewan}. Di situ ia
menemukan makanannya, air, perlindungan, dan tempat yang aman bagi
anaknya.
\end{solution}

\begin{solution}{exo:g2:where-animals-live:2}
Kolam: katak, capung (atau ikan duri, siput kolam). Hutan: tupai,
burung pelatuk (atau rusa, semut kayu). Padang rumput: kelinci,
belalang (atau kupu-kupu, tikus ladang).
\end{solution}

\begin{solution}{exo:g2:where-animals-live:3}
Ikan duri: kolamnya. Burung pelatuk: hutannya. Belalang: padang
rumputnya. Kutu kayu: di bawah batu sebuah tembok tua (sudut lembap
yang tersembunyi).
\end{solution}

\begin{solution}{exo:g2:where-animals-live:4}
Kataknya: kaki berselaput yang mendorong air. Tupainya: cakar
pencengkeram dan ekor besar sebagai penyeimbang.
\end{solution}

\begin{solution}{exo:g2:where-animals-live:5}
Tangan sekop penggali yang lebar; \omterm{def:g1:animals-around-us:covering}{bulu} yang terbaring rata ke kedua
arah sehingga tanah tidak pernah mengacaknya; \omterm{def:g1:the-five-senses:sense}{mata} yang nyaris tidak
dibutuhkannya dalam gelap.
\end{solution}

\begin{solution}{exo:g2:where-animals-live:6}
Burung pelatuk: tinggi di \omterm{def:g1:plants-around-us:parts}{batang} pohon. Tupai: di dahan. Rusa: di
antara semak. Semut kayu: di tanah, dalam sarang gundukannya.
\end{solution}

\begin{solution}{exo:g2:where-animals-live:7}
Mereka lenyap bersama kolamnya --- tanpa \omterm{def:g2:where-animals-live:habitat}{habitatnya} mereka kehilangan
makanan, perlindungan dan tempat berkembang biak sekaligus. Jadi
melindungi \omterm{def:g1:animals-around-us:animal}{hewan} bermula dengan melindungi tempat mereka hidup.
\end{solution}

\begin{solution}{exo:g2:where-animals-live:8}
Jawaban terbuka. Bahkan satu batu yang dibalikkan sering menaungi kutu
kayu, semut, seekor kumbang, kadang seekor kelabang --- \omterm{def:g2:where-animals-live:habitat}{habitat} kecil
ternyata sangat penuh.
\end{solution}

\begin{solution}{exo:g2:where-animals-live:9}
Kaki berselaput (mengayuh), \omterm{def:g1:animals-around-us:covering}{bulu} yang menolak air, paruh penyaring:
ketiganya menunjuk ke air --- \omterm{def:g2:where-animals-live:habitat}{habitat} kolam atau danau. \omterm{def:g1:animals-around-us:animal}{Hewannya}
sangat mungkin seekor bebek.
\end{solution}

\begin{solution}{exo:g2:where-animals-live:10}
Perlengkapan kataknya --- kaki berselaput, \omterm{def:g1:the-five-senses:sense}{kulit} yang lembap, tubuh
yang dibuat untuk air dan tepian yang lembap --- tidak berguna di
dalam kotak kering, dan \omterm{def:g2:where-animals-live:habitat}{habitatnya} memberinya lebih daripada makanan:
air bagi \omterm{def:g1:the-five-senses:sense}{kulit} dan \omterm{def:g2:animals-grow-up:born}{telurnya}, tempat bersembunyi, katak-katak lain.
Berbaik hati kepada \omterm{def:g1:animals-around-us:animal}{hewan} liar berarti membiarkannya di dalam --- atau
mengembalikannya ke --- \omterm{def:g2:where-animals-live:habitat}{habitatnya}.
\end{solution}
